% !TEX root = ../main.tex
%----------------------------------------------------------
% 00-abstract.tex — Abstract
%----------------------------------------------------------

\begin{abstract}
Possibility theory offers a natural framework for encoding epistemic
uncertainty in hazard forecasting, yet possibilistic predictions lack
a dedicated verification methodology.  A subnormal possibility
distribution---one whose maximum falls below unity---decomposes the
forecast signal into three components: possibility $\poss$ (how
plausible is each outcome?), ignorance $\ign$ (how much of the
plausibility budget is unassigned?), and conditional necessity $\Nc$
(among covered outcomes, which dominates?).  The ignorance signal has
no probabilistic counterpart: probability, which must sum to unity,
cannot distinguish confident forecasts from uncertain ones.
This paper presents a \emph{five-number scorecard} that evaluates
these three components once the outcome is known, comprising
depth-of-truth ($\alpha^*$), diffuseness ($\eta$), support margin
($\delta$), ignorance ($\ign$), and conditional necessity of truth
($\Nc^*$).  A possibility-to-probability conversion preserves
ignorance as an explicit outcome, enabling information-gain
decomposition.  Together, the scorecard and the derived probabilities
define three verification facets---categorical, probabilistic, and
native possibilistic---each diagnosing failure modes that may be
obscured in any single facet.  All concepts are illustrated with Storm
Prediction Center convective outlook categories (\spc{NONE} through
\spc{HIGH}).
\end{abstract}

\keywords{possibility theory, forecast verification, subnormal distributions, scoring rules, severe weather}
